\paragrafos{ΕΚΠ και ΜΚΔ: πότε ξανασυναντιόμαστε, πόσο μεγάλα κομμάτια}{Α.1.4, «ΕΚΠ και ΜΚΔ φυσικών αριθμών»}

\begin{anagke}
Η γιαγιά Μυρτώ ποτίζει τα φυτά της αυλής κάθε 3 ημέρες και ο παππούς Γιάννης κλαδεύει τη λεμονιά κάθε 4 ημέρες. Σήμερα έκαναν και οι δύο τη δουλειά τους. «Πότε θα τα ξανακάνουμε και οι δύο την ίδια μέρα;» ρωτά η Ζωή, που έχει έρθει για επίσκεψη. Ο Ηλίας γράφει στο τετράδιο τις ημέρες της γιαγιάς: 3, 6, 9, 12, 15, 18, και του παππού: 4, 8, 12, 16, 20. «Στις 12!» φωνάζει.

Το απόγευμα η θεία Μαρία φέρνει 48 μανταρίνια και 60 πορτοκάλια και θέλει να τα βάλει σε ίσα καλάθια, με τον ίδιο αριθμό μανταρινιών και πορτοκαλιών σε κάθε καλάθι, χωρίς να περισσέψει τίποτα. «Πόσα το πολύ καλάθια μπορώ να φτιάξω;» ρωτά. Αυτό δεν είναι «πότε ξανασυναντιόμαστε»: είναι «σε πόσα ίσα κομμάτια χωρίζονται και τα δύο».

Δύο διαφορετικά προβλήματα: το ένα ψάχνει έναν αριθμό που τον φτάνουν και οι δύο (πολλαπλάσιο), το άλλο έναν αριθμό που χωράει σε όλους (διαιρέτη). Σε αυτή την παράγραφο μαθαίνουμε το ελάχιστο κοινό πολλαπλάσιο (ΕΚΠ) και τον μέγιστο κοινό διαιρέτη (ΜΚΔ), και πώς τα βρίσκουμε με την ανάλυση σε γινόμενο πρώτων παραγόντων και με τον αλγόριθμο του Ευκλείδη.
\end{anagke}

\begin{lexi}{πολλαπλάσιο, κοινό, ελάχιστο, μέγιστο}{πολλαπλάσια, ΕΚΠ, ΜΚΔ}
Το \lexiword{πολλαπλάσιο} ενός αριθμού είναι αυτός που προκύπτει όταν τον πάρουμε πολλές φορές, δηλαδή τον πολλαπλασιάσουμε. \lexiword{Κοινό} είναι αυτό που ανήκει και στους δύο. Με το ΕΚΠ ψάχνουμε το \lexiword{ελάχιστο}, το πιο μικρό, κοινό πολλαπλάσιο, και με το ΜΚΔ τον \lexiword{μέγιστο}, τον πιο μεγάλο, κοινό διαιρέτη. Το $0$ είναι πολλαπλάσιο κάθε αριθμού, γι' αυτό στο ΕΚΠ ζητάμε το μικρότερο \emph{μη μηδενικό} κοινό πολλαπλάσιο.
\end{lexi}

\begin{metafrasi}
Οι αριθμοί που προκύπτουν όταν πολλαπλασιάσω το 4. & $0,4,8,12,16,\dots$ (πολλαπλάσια του $4$)\\
Το πρώτο κοινό πολλαπλάσιο των 4 και 6. & $\EKP(4,6)=12$\\
Το μεγαλύτερο κοινό τους κομμάτι των 12 και 18. & $\MKD(12,18)=6$\\
Δύο αριθμοί δεν έχουν κοινό διαιρέτη εκτός από το 1. & $\MKD(\alpha,\beta)=1$, πρώτοι μεταξύ τους\\
\end{metafrasi}

\tmima{Η θεωρία}

\tmimasmall{Πολλαπλάσια και ΕΚΠ}

\begin{orismos}{πολλαπλάσια}
\textbf{Πολλαπλάσια} ενός φυσικού $\alpha$ είναι οι αριθμοί που προκύπτουν όταν τον πολλαπλασιάσουμε με όλους τους φυσικούς: $0,\ \alpha,\ 2\alpha,\ 3\alpha,\ 4\alpha,\dots$ Για παράδειγμα τα πολλαπλάσια του $9$ είναι τα $0,9,18,27,36,45,\dots$
\end{orismos}

\begin{orismos}{ελάχιστο κοινό πολλαπλάσιο}
Το \textbf{ελάχιστο κοινό πολλαπλάσιο} (ΕΚΠ) δύο ή περισσότερων φυσικών αριθμών που δεν είναι μηδέν είναι το μικρότερο (μη μηδενικό) από τα κοινά πολλαπλάσια των αριθμών αυτών.
\end{orismos}

Παράδειγμα: τα πολλαπλάσια του $4$ είναι $0,4,8,\mathbf{12},16,20,\mathbf{24},28,32,\mathbf{36},\dots$ και του $6$ είναι $0,6,\mathbf{12},18,\mathbf{24},30,\mathbf{36},42,\dots$ Τα κοινά είναι τα $12,24,36,\dots$ (άπειρα) και το μικρότερο είναι το $12$, άρα $\EKP(4,6)=12$.

\tmimasmall{Διαιρέτες και ΜΚΔ}

Από την προηγούμενη παράγραφο ξέρουμε ότι \textbf{διαιρέτες} ενός αριθμού είναι όσοι τον διαιρούν τέλεια. Ισχύουν τα εξής:
\begin{itemize}[leftmargin=1.2em,itemsep=0.2ex,topsep=2pt]
\item Κάθε φυσικός (εκτός από το $0$) διαιρεί τα πολλαπλάσιά του.
\item Κάθε φυσικός που διαιρείται από έναν άλλον είναι πολλαπλάσιο αυτού του αριθμού.
\item Αν ένας φυσικός διαιρεί έναν άλλον, θα διαιρεί και τα πολλαπλάσια του άλλου.
\end{itemize}

\begin{orismos}{μέγιστος κοινός διαιρέτης}
Ο \textbf{μέγιστος κοινός διαιρέτης} (ΜΚΔ) δύο ή περισσότερων φυσικών αριθμών είναι ο μεγαλύτερος από τους κοινούς διαιρέτες τους.
\end{orismos}

Παράδειγμα: οι διαιρέτες του $12$ είναι οι $1,2,3,4,6,12$ και του $18$ οι $1,2,3,6,9,18$. Οι κοινοί διαιρέτες είναι οι $1,2,3,6$ και ο μεγαλύτερος είναι ο $6$, άρα $\MKD(12,18)=6$.

Αν $\MKD(\alpha,\beta)=1$, οι αριθμοί $\alpha$ και $\beta$ λέγονται \textbf{πρώτοι μεταξύ τους}. Για παράδειγμα οι $6$ και $7$ έχουν κοινό διαιρέτη μόνο το $1$, άρα είναι πρώτοι μεταξύ τους (αν και ο $6$ δεν είναι πρώτος αριθμός).

\begin{prosoxi}
Το «πρώτοι μεταξύ τους» δεν σημαίνει ότι οι δύο αριθμοί είναι πρώτοι. Οι $8$ και $15$ δεν είναι πρώτοι αριθμοί, αλλά είναι πρώτοι μεταξύ τους.
\end{prosoxi}

\tmimasmall{ΕΚΠ και ΜΚΔ με ανάλυση σε πρώτους παράγοντες}

Η λίστα με τα πολλαπλάσια ή τους διαιρέτες είναι κουραστική για μεγάλους αριθμούς. Αναλύουμε πρώτα όλους τους αριθμούς σε γινόμενο πρώτων παραγόντων.

\begin{kanonas}{ΕΚΠ και ΜΚΔ με ανάλυση}
Το \textbf{ΕΚΠ} είναι το γινόμενο όλων των πρώτων παραγόντων (κοινών και μη κοινών) με τον \textbf{μεγαλύτερο} εκθέτη τους.\\
Ο \textbf{ΜΚΔ} είναι το γινόμενο μόνο των \textbf{κοινών} πρώτων παραγόντων με τον \textbf{μικρότερο} εκθέτη τους.
\end{kanonas}

\tmimasmall{Ο αλγόριθμος του Ευκλείδη για το ΜΚΔ}

Ο Ευκλείδης έδωσε μια μέθοδο που βρίσκει τον ΜΚΔ χωρίς καμία ανάλυση. Τα βήματα, για τους αριθμούς $125,170,280$:

\begin{center}
\renewcommand{\arraystretch}{1.3}\setlength{\tabcolsep}{9pt}
\begin{tabular}{ccc}
$125$ & $170$ & $280$\\
$125$ & $45$ & $30$\\
$5$ & $15$ & $30$\\
$5$ & $0$ & $0$
\end{tabular}
\end{center}

Γράφουμε τους αριθμούς και βρίσκουμε τον μικρότερο. Στην επόμενη γραμμή γράφουμε \emph{τον μικρότερο αριθμό ως έχει} και κάτω από κάθε άλλον το \emph{υπόλοιπο της διαίρεσής του με τον μικρότερο} ($170=125\cdot 1+45$, $280=125\cdot 2+30$). Επαναλαμβάνουμε με τον νέο μικρότερο μη μηδενικό αριθμό. Σταματάμε όταν όλοι οι άλλοι γίνουν $0$. Ο ΜΚΔ είναι ο μη μηδενικός αριθμός της τελευταίας γραμμής: $\MKD(125,170,280)=5$.

\tmima{Λυμένα παραδείγματα}

\paradeigma{Το λεωφορείο και το τρένο}
\tixero Από το χωριό φεύγει λεωφορείο κάθε $12$ λεπτά και τρένο κάθε $15$ λεπτά. Έφυγαν μαζί στις $8{:}00$.\par
\tithelo Πότε θα ξαναφύγουν μαζί;\par
\possindeo Ζητάω το μικρότερο κοινό πολλαπλάσιο των $12$ και $15$. Πολλαπλάσια του $12$: $12,24,36,48,60,\dots$ Πολλαπλάσια του $15$: $15,30,45,60,\dots$ Το πρώτο κοινό είναι το $60$.\par
\apantisi $\EKP(12,15)=60$, άρα μετά από $60$ λεπτά, στις $9{:}00$.\par
\elegxos Το $60=12\cdot 5=15\cdot 4$.

\paradeigma{Οι χάντρες της Αγγελικής}
\tixero Η Αγγελική έχει $36$ κόκκινες, $48$ κίτρινες και $60$ μπλε χάντρες. Φτιάχνει όσο το δυνατόν περισσότερα ίδια κολιέ, με τον ίδιο αριθμό χαντρών κάθε χρώματος, χωρίς να περισσέψει καμία.\par
\tithelo Πόσα κολιέ φτιάχνει και τι έχει το καθένα;\par
\possindeo Ο αριθμός των κολιέ πρέπει να διαιρεί και τα τρία πλήθη, και ζητάω τον μεγαλύτερο, άρα τον ΜΚΔ. Αναλύω: $36=2^2\cdot 3^2$, $48=2^4\cdot 3$, $60=2^2\cdot 3\cdot 5$. Κοινοί παράγοντες: $2$ και $3$ με μικρότερους εκθέτες $2$ και $1$. Άρα $\MKD=2^2\cdot 3=12$. Κάθε κολιέ έχει $36:12=3$ κόκκινες, $48:12=4$ κίτρινες και $60:12=5$ μπλε.\par
\apantisi $12$ κολιέ, με $3$ κόκκινες, $4$ κίτρινες και $5$ μπλε χάντρες.\par
\elegxos $12\cdot 3=36$, $12\cdot 4=48$, $12\cdot 5=60$.

\paradeigma{ΕΚΠ και ΜΚΔ με ανάλυση, και με Ευκλείδη}
\tixero Οι αριθμοί $84$, $90$ και $126$.\par
\tithelo Το ΕΚΠ και ο ΜΚΔ τους.\par
\possindeo Αναλύω: $84=2^2\cdot 3\cdot 7$, $90=2\cdot 3^2\cdot 5$, $126=2\cdot 3^2\cdot 7$. Όλοι οι πρώτοι παράγοντες είναι οι $2,3,5,7$, με μεγαλύτερους εκθέτες $2,2,1,1$: $\EKP=2^2\cdot 3^2\cdot 5\cdot 7=1.260$. Κοινοί παράγοντες είναι μόνο οι $2$ και $3$, με μικρότερους εκθέτες $1,1$: $\MKD=2\cdot 3=6$. Με τον αλγόριθμο του Ευκλείδη: $(84,90,126)\to(84,6,42)\to(0,6,0)$, άρα ΜΚΔ $=6$.\par
\apantisi $\EKP=1.260$ και $\MKD=6$.\par
\elegxos Το $1.260:84=15$, $1.260:90=14$, $1.260:126=10$ και ο $6$ διαιρεί τους $84,90,126$.

\tmima{Ασκήσεις}

\zoni{Το κατάλαβα}{ορισμοί και υπολογισμοί}

\ex Συμπλήρωσε: α) Αν $\MKD(\alpha,\beta)=1$, οι $\alpha$ και $\beta$ λέγονται \dots μεταξύ τους. β) Κάθε φυσικός \dots τα πολλαπλάσιά του. γ) Κάθε φυσικός που διαιρείται από έναν άλλον είναι \dots αυτού του αριθμού. δ) Ο ΜΚΔ δύο αριθμών είναι ο \dots από τους κοινούς διαιρέτες τους.

\ex Γράψε τα $8$ πρώτα πολλαπλάσια των $7$ και $12$ (ξεκινώντας από το $0$).

\ex Γράψε τα $8$ πρώτα μη μηδενικά πολλαπλάσια των $6$ και $9$. Ποια είναι κοινά; Ποιο είναι το $\EKP(6,9)$;

\ex Γράψε όλους τους διαιρέτες των $18$, $30$ και $36$. Βρες τους κοινούς διαιρέτες και τον $\MKD(18,30,36)$.

\ex Είναι πρώτοι μεταξύ τους οι αριθμοί; α) $8$ και $15$ \quad β) $14$ και $21$ \quad γ) $22$ και $45$ \quad δ) $9$ και $28$ \quad ε) $12$ και $35$.

\ex Συμπλήρωσε τον πίνακα με τον ΜΚΔ και το ΕΚΠ κάθε ζεύγους: $(6,8)$, $(10,15)$, $(4,9)$, $(14,21)$, $(16,24)$.

\ex Ανάλυσε τους $36$ και $54$ σε γινόμενο πρώτων παραγόντων και βρες τον ΜΚΔ και το ΕΚΠ τους.

\zoni{Το χρησιμοποιώ}{ΕΚΠ ή ΜΚΔ;}

\ex \textbf{Οι προπονήσεις.} Ο Νίκος προπονείται κάθε $3$ ημέρες, ο Βασίλης κάθε $5$ και η Πολυξένη κάθε $6$. Σήμερα προπονήθηκαν όλοι μαζί. Μετά από πόσες ημέρες θα ξανασυναντηθούν;

\ex \textbf{Τα καλάθια.} Η θεία Μαρία έχει $54$ μήλα και $72$ αχλάδια και τα βάζει σε ίσα καλάθια, χωρίς να περισσέψει τίποτα και με τον ίδιο αριθμό μήλων και αχλαδιών σε κάθε καλάθι. Πόσα το πολύ καλάθια μπορεί να φτιάξει και τι έχει το καθένα;

\ex \textbf{Τα πλακάκια.} Το δωμάτιο έχει διαστάσεις $240\,\mathrm{cm}\times 180\,\mathrm{cm}$ και θέλουμε να το στρώσουμε με τετράγωνα πλακάκια, ίδια, χωρίς κοψίματα. Ποια είναι η μεγαλύτερη δυνατή πλευρά του πλακιδίου και πόσα πλακάκια χρειάζονται;

\ex \textbf{Τα φωτάκια.} Στη γιορτή του χωριού ένα κόκκινο φωτάκι ανάβει κάθε $12$ δευτερόλεπτα και ένα πράσινο κάθε $18$. Άναψαν μαζί. Μετά από πόσα δευτερόλεπτα θα ανάψουν ξανά μαζί;

\ex \textbf{Οι σελιδοδείκτες.} Η Φανή έχει $40$ σελιδοδείκτες και $56$ αυτοκόλλητα. Φτιάχνει ίδια πακέτα, με τον ίδιο αριθμό από κάθε είδος, χωρίς να περισσέψει τίποτα. Πόσα το πολύ πακέτα φτιάχνει και τι έχει κάθε πακέτο;

\ex \textbf{Ο φούρνος.} Ένας φούρνος έχει $60$ κουλούρια, $84$ κρουασάν και $108$ τσουρέκια και τα βάζει σε ίδια κουτιά, με τον ίδιο αριθμό από κάθε είδος και χωρίς να περισσέψει τίποτα. Πόσα το πολύ κουτιά γεμίζει και τι έχει το καθένα;

\zoni{Το σκέφτομαι}{ερωτήσεις με λίγη περισσότερη σκέψη}

\ex \textbf{Μια ανακάλυψη.} Για τα ζεύγη $(4,6)$, $(12,18)$ και $(10,15)$ βρες τον ΜΚΔ και το ΕΚΠ. Μετά υπολόγισε το γινόμενο $\MKD\cdot\EKP$ και το γινόμενο των δύο αριθμών κάθε ζεύγους. Τι παρατηρείς;

\ex Βρες δύο διαφορετικούς φυσικούς που έχουν $\MKD=6$ και $\EKP=60$.

\ex α) Βρες το $\EKP(8,15)$ και το $\EKP(9,10)$. Πώς σχετίζονται με το γινόμενο των αριθμών, όταν οι αριθμοί είναι πρώτοι μεταξύ τους; β) Βρες τον $\MKD$ και το $\EKP$ των $6$ και $30$ (ο ένας διαιρεί τον άλλον). Τι παρατηρείς;

\zoni{Για το σπίτι}{εξάσκηση}

\ex Βρες με ανάλυση τον ΜΚΔ και το ΕΚΠ των $325$ και $338$ (από την παράγραφο 1.3) και των $24$, $42$, $54$.

\ex Βρες με τον αλγόριθμο του Ευκλείδη τον ΜΚΔ των $48$, $80$, $120$. Μετά εξέτασε αν οι $1.000$ και $1.303$ είναι πρώτοι μεταξύ τους.
